\section{The structure of an FsLex specification}
A scanner specification has a header containing F\# imports and helpers, named regular expressions,
and one or more scanner rules. Unlike dynamically adding attributes to a token object, F\# makes each token's payload
part of a discriminated union. A number token carries an integer, and a symbol token carries its character.
The arithmetic notebook uses the following shared types.
@@snippet:scanner-types@@

A rule has the form \texttt{pattern \{ F\# action \}}. The action may return a token or call the scanner again
to discard whitespace. FsLex selects the longest matching input prefix; among equally long matches the earlier rule wins.
The following complete specification recognizes natural numbers, arithmetic operators and parentheses.
It preserves an explicit illegal-character token, rather than silently dropping unexpected input.
@@snippet:scanner-arithmetic@@

The newline action advances \texttt{EndPos}. This is separate from returning a token, so locations remain available
even when whitespace is skipped. EOF is an explicit token that ends iteration. The driver constructs a fresh
buffer, sets its initial position and collects tokens until that marker.
@@snippet:scanner-driver@@
\noindent Execute \href{../../Chapter-03/01-Ply-Scanning-Example.ipynb}{\nolinkurl{Chapter-03/01-Ply-Scanning-Example.ipynb}}.
Its filename retains the source notebook's name for direct comparison; its scanner is implemented with FsLex.

\section{Practical regular expressions in .NET}
The lecture's algebraic regular expressions use alternatives, concatenation and Kleene star.
The .NET library adds character classes, capturing groups, assertions, bounded repetition and matching options.
These operations are library calls, separate from FsLex's scanner-specification syntax.
F\# verbatim strings begin with \texttt{@} so that regex backslashes are written directly.
For example, \texttt{@"[0-9]+"} matches a positive-length digit sequence.
The tutorial defines reusable matching helpers as follows.
@@snippet:regex-api@@

\texttt{Regex.Matches} returns match objects; their \texttt{Value} is the matched text and
\texttt{Groups} stores captures. \texttt{RegexOptions.IgnoreCase} and \texttt{RegexOptions.Multiline}
control matching explicitly. The \texttt{\textbackslash G} anchor is useful for complete tokenization:
it requires the next match to begin at the previous match's end, so unmatched punctuation cannot disappear.
\href{../../Chapter-03/02-Regexp-Tutorial.ipynb}{\nolinkurl{Chapter-03/02-Regexp-Tutorial.ipynb}}
retains the original searches, grouping examples, lookarounds and analysis of \textsl{Alice in Wonderland}.

\section{Evaluating an examination}
The examination format contains headers, a maximum score, student identifiers and individual exercise scores.
The scanner returns typed tokens instead of keeping evaluation state inside the lexer.
A separate evaluator groups scores by student and computes the final grade. This boundary lets us inspect the token stream
before deciding whether a grading error originated in scanning or calculation.
@@snippet:exam-scanner@@
@@snippet:exam-grade@@

See \href{../../Chapter-03/03-Exam-Evaluation.ipynb}{\nolinkurl{Chapter-03/03-Exam-Evaluation.ipynb}}
and its file-based counterpart. Their data, grade calculations and chart examples are executable notebook cells.

\section{Scanner states}
A lexical mode selects which patterns are currently active. FsLex expresses modes naturally as mutually recursive
scanner rules. The HTML example enters a \texttt{header} or \texttt{script} rule and returns to the main rule
when the corresponding closing tag is found. EOF inside an unfinished mode is reported as an error.
No runtime state-name registry is needed: the active F\# function represents the mode.
@@snippet:html-states@@

\texttt{WebUtility.HtmlDecode} handles named and numeric entities through .NET.
This scanner is an exercise over the lecture's constrained HTML format, not a full HTML5 tree builder.
\href{../../Chapter-03/04-Html2Text.ipynb}{\nolinkurl{Chapter-03/04-Html2Text.ipynb}}
and \href{../../Chapter-03/05-LaTeX2HTML/LaTeX2HTML.ipynb}{\nolinkurl{Chapter-03/05-LaTeX2HTML/LaTeX2HTML.ipynb}}
extend the same separation between lexical modes, text decoding and output formatting.

\section{Check your understanding}
\begin{enumerate}
\item Explain longest-match selection and what happens when two rules match equally long prefixes.
\item Why are the input alphabet, an epsilon transition and an EOF token different concepts?
\item Modify the arithmetic scanner to include identifier tokens with string payloads.
\item Explain how the token type and grading evaluator prevent a student identifier from becoming a score.
\item Add a lexical mode for comments and report an unfinished comment at EOF.
\item Compare the expressive power of the lecture's regular expressions with .NET backreferences.
\end{enumerate}
